\subsection{Profil Desktop-cloud}
\label{subsec:profil-desktop-cloud}

\texttt{desktop-cloud} aktiviert \texttt{frontend}, \texttt{backend},
\texttt{tauri} und \texttt{cloud}. Es ergänzt \texttt{desktop-local} damit um
FastAPI und die Deployment-Grenze: Das gemeinsame Frontend wird in der
Tauri-Webview ausgeführt und greift über die konfigurierte öffentliche URL auf
das getrennte Backend zu. Datenbankinfrastruktur bleibt ohne Capability
ausgeschlossen; PostgreSQL kann wegen des vorhandenen Backends ergänzt werden
\parencite{kleiveist2026profiles,kleiveist2026deployment}.

Die Basisfeatures und Feature-Pfade sind identisch zu
\texttt{full-platform}. Die Profil-ID dokumentiert jedoch den Desktop-Client
als vorgesehenen Produktkanal, während das folgende Profil Browser und Desktop
gemeinsam adressiert. Ein zusätzlicher technischer Baustein entsteht aus
dieser semantischen Abgrenzung nicht.
